\documentclass[10pt,a4paper]{amsart}
\usepackage[margin=19mm]{geometry}
\usepackage{amsmath,amssymb,mathtools}
\usepackage[scr=rsfs,cal=euler]{mathalfa}
\usepackage{xfrac}
\usepackage[unicode,hidelinks,bookmarks=false]{hyperref}
\newcommand{\ie}{i.e.\ }
\newcommand{\cf}{cf.\ }
\newcommand{\resp}{respectively\ }
\newcommand{\Subsection}{\subsection}
\providecommand{\nobreakdash}{\nobreak\hbox{-}}
\setlength{\emergencystretch}{2em}
\multlinegap0pt
\usepackage{amssymb}
\usepackage{mathtools}
\newtheorem{lemma}{Lemma}
\newtheorem{proposition}{Proposition}
\newcommand{\Jp}{\mathcal J_p}
\newcommand{\R}{\mathbb R}
\newcommand{\dd}{\,d}
\title{Gradient H\"older regularity for singular fractional $p$-Laplace equations}
\author{Chao Zhang}
\date{Source: arXiv:2608.16243v2 (CC BY 4.0)}
\begin{document}
\maketitle

\noindent\emph{Source and licence.} Gradient H\"older regularity for singular fractional $p$-Laplace equations. Version arXiv:2608.16243v2 (CC BY 4.0), distributed by arXiv under the Creative Commons Attribution 4.0 International licence, http://creativecommons.org/licenses/by/4.0/, as recorded in the arXiv licence metadata for this exact version and shown on the abstract page https://arxiv.org/abs/2608.16243v2. Full licence notice: \texttt{licenses/arxiv-2608.16243-CC-BY-4.0.txt}.\par\smallskip

\noindent\textit{Editorial scope.} Counted unit: the article's proposition prop:coarea-weak-linearization (source0.tex lines 3881--3933). The article's complete original proof of it is delivered with it (source0.tex lines 3935--4257, one \texttt{proof} environment of 323 lines). No mathematics is written, repaired or re-derived: the statement and the proof are the article's own lines, in the article's own notation, and no retained line is substituted or re-flowed.  Retained uncounted as the source-local context the proof needs: lines 273--278; lines 3062--3102; lines 3643--3650.  Nothing was omitted from any retained span, so the package carries no omission record.  No retained line contains a citation, so the wrapper carries no bibliography and no bibliography entry.  No retained line contains a figure environment, an image-inclusion command or drawing code, so no image is delivered and the package declares no asset dependency.  Editorial formatting only, in addition: an AMS \texttt{amsart} preamble that restates the article's own declarations and macros used by the retained lines, namely the environments \texttt{lemma}, \texttt{proposition} on separate counters, together with the standard packages those lines need.  The retained environments print in this document as Lemma 1, Proposition 1 in order of appearance, whereas the article prints the same items with the numbering of its own full document.  The complete native source of the article as distributed in the arXiv e-print for this exact version is bundled unchanged as \texttt{sources/207816/source0.tex}.
\begin{equation}
 \operatorname{P.V.}\int_{\R^n}
 \frac{\Jp(U(x)-U(y))}{|x-y|^{n+sp}}\dd y=0
 \qquad\text{in }\R^n.
 \label{eq:equation}
\end{equation}

\begin{lemma}[Regional angular compactness at the gradient contact]
\label{lem:regional-angular-compactness}
Let $1<p<2$, $0<\alpha<1$, and let $U\in C^{1,\alpha}(B_{2r_0})$ satisfy
\begin{equation}
 \|\nabla U-q\|_{L^\infty(B_{2r_0})}\le\frac14|q|,
 \qquad q\ne0.
 \label{eq:noncritical-gradient-tube}
\end{equation}
Fix a unit vector $\xi$.  If $x$, $x-\rho\omega$ and their translates by
$h\xi$ belong to $B_{2r_0}$, put
\begin{align*}
 A_0(x,\rho,\omega)&:=U(x)-U(x-\rho\omega),\\
 A_h(x,\rho,\omega)&:=U(x+h\xi)-U(x-\rho\omega+h\xi),
\end{align*}
and
\begin{equation*}
 b_h(x,\rho,\omega)
 :=(p-1)\rho^{2-p}\int_0^1
 |(1-t)A_0+tA_h|^{p-2}\dd t.
\end{equation*}
Then, uniformly for the admissible $x$, $\rho$ and $h$,
\begin{equation}
 c\le b_h(x,\rho,\omega)\quad\text{for a.e. }\omega,
 \qquad
 \int_{\mathbb S^{n-1}}b_h(x,\rho,\omega)\dd\omega\le C.
 \label{eq:regional-angular-ellipticity}
\end{equation}
Moreover, if $|h|\le\rho$, then
\begin{equation}
 \int_{\mathbb S^{n-1}}
 \left|b_h(x,\rho,\omega)
 -(p-1)\left|\frac{A_0(x,\rho,\omega)}{\rho}\right|^{p-2}
 \right|\dd\omega
 \le C|h|^{\alpha(p-1)}.
 \label{eq:regional-angular-convergence}
\end{equation}
Here the constants depend only on $n,p,|q|$ and
$[\nabla U]_{C^\alpha(B_{2r_0})}$, after the geometry has been normalized
so that the displayed segments stay a fixed positive distance from the
boundary.
\end{lemma}

\begin{equation}
 \int A\psi^2\dd x
 \le \epsilon\mathcal E_K(\psi,\psi)
 +C_*(1+M)^{\tau/(\tau-\kappa)}
       \epsilon^{-\kappa/(\tau-\kappa)}\|\psi\|_2^2,
 \qquad 0<\epsilon\le1,
 \label{eq:kato-infinitesimal-form-bound}
\end{equation}

\begin{proposition}[Directional coarea and the full weak linearization]
\label{prop:coarea-weak-linearization}
Let $U$ be a globally $L$-Lipschitz entire solution of
\eqref{eq:equation}.  Suppose that, for some $q\ne0$ and
$e=q/|q|$,
\begin{equation}
 U\in C^{1,\alpha}(B_{2r_0}),
 \qquad
 \|\nabla U-q\|_{L^\infty(B_{2r_0})}\le\frac14|q|.
 \label{eq:strict-directional-monotonicity}
\end{equation}
In particular, $\partial_eU\ge\lambda_0:=3|q|/4$ in $B_{2r_0}$.
Set
\[
 \tau:=sp-p+2,
 \qquad
 J_U(x,y):=(p-1)
 \frac{|U(x)-U(y)|^{p-2}}{|x-y|^{n+sp}}.
\]
On the set \(U(x)=U(y)\) we set \(J_U(x,y)=0\).  For every fixed
\(y\), this set has zero \(x\)-measure locally in \(B_{2r_0}\), by
strict monotonicity on the \(e\)-lines, so the convention does not
change any of the forms below.
Then $J_U$ is locally a well-defined symmetric Dirichlet kernel.  If
$B_r(x_0)\subset B_{r_0/2}$ and
\[
 A_U(x):=\int_{\R^n\setminus B_{r_0}}J_U(x,y)\dd y,
\]
then
\begin{equation}
 \int_{B_r(x_0)}A_U(x)\dd x\le Cr^{n+p-2}.
 \label{eq:crossed-morrey-bound}
\end{equation}
More generally, for every \(D_0\Subset B_{r_0}\) the same estimate
holds for balls \(B_r(x_0)\subset D_0\), with a constant which may also
depend on \(\operatorname{dist}(D_0,\partial B_{r_0})\).
In particular, $A_U$ is a Morrey--Kato absorption for the regional order
$\tau$, since
\begin{equation}
 \tau-(2-p)=sp>0.
 \label{eq:strict-kato-gap}
\end{equation}

Moreover, the a.e. directional derivative \(w=\partial_eU\) belongs to
\(\mathcal F_{J_U,\mathrm{loc}}(B_{r_0})\) and satisfies
\begin{equation}
 \iint_{\R^n\times\R^n}
 J_U(x,y)\delta w(x,y)\delta\varphi(x,y)\dd x\!\dd y=0
\label{eq:full-weak-linearization}
\end{equation}
for every \(\varphi\in C_c^\infty(B_{r_0})\).  The identity extends to
every \(\varphi\in\mathcal F_{J_U,c}(B_{r_0})\).
\end{proposition}

\begin{proof}
We first record the quantitative form of the directional coarea
estimate used below.  For \(t\in[0,1]\) and \(0\le h<r_0/4\), put
\[
 W_{t,h}(x):=(1-t)U(x)+tU(x+he),
\]
with \(W_{t,0}=U\).  If \(K\Subset B_{r_0}\), then, after decreasing
the upper bound for \(h\) by a number depending only on \(K\),
\begin{equation}
 \sup_{\substack{0\le h<h_K\\t\in[0,1]\\a\in\R}}
 \int_E|W_{t,h}(x)-a|^{p-2}\dd x
 \le C_K|E|^{p-1}
 \qquad(E\subset K).
 \label{eq:quantitative-directional-coarea}
\end{equation}
Indeed, write \(x=x'+\ell e\), with \(x'\in e^\perp\).
On every line section, \(\partial_eW_{t,h}\ge\lambda_0\).  The
one-dimensional area formula, followed by decreasing rearrangement
about \(a\), gives
\[
 \int_{E_{x'}}|W_{t,h}(x'+\ell e)-a|^{p-2}\dd\ell
 \le C|E_{x'}|^{p-1}.
\]
Integration in \(x'\), followed by H\"older's inequality on the
projection of \(K\), proves
\eqref{eq:quantitative-directional-coarea}.

For \(E=B_r(x_0)\), the line calculation gives the sharper bound
\(Cr^{n-1}r^{p-1}=Cr^{n+p-2}\).  Apply it with \(h=0\) and
\(a=U(y)\), and then integrate in \(y\notin B_{r_0}\).  The two sets
are separated, while for large \(y\) the remaining kernel is bounded
by \(C(1+|y|)^{-n-sp}\), uniformly in \(x\in B_r(x_0)\).
This proves \eqref{eq:crossed-morrey-bound}.
Replacing \(B_{r_0/2}\) by any fixed \(D_0\Subset B_{r_0}\) in the
same argument proves the stated local version, since the two regions
remain positively separated.

We next justify differentiation of the equation.  For $0<h<r_0/4$ put
\begin{align*}
 U_h(x)&:=U(x+he),\qquad w_h:=\frac{U_h-U}{h},\\
 c_h(x,y)&:=(p-1)\int_0^1
 |(1-t)\delta U(x,y)+t\delta U_h(x,y)|^{p-2}\dd t.
\end{align*}
Here and below \(c_h=0\) when
\(\delta U(x,y)=\delta U_h(x,y)=0\).  The integral is finite otherwise,
even if the affine segment between the two increments crosses zero,
because \(p-2>-1\).
Subtraction of the two entire weak equations gives the exact identity
\begin{equation}
 \iint_{\R^n\times\R^n}
 \frac{c_h(x,y)\delta w_h(x,y)\delta\varphi(x,y)}
 {|x-y|^{n+sp}}\dd x\!\dd y=0.
 \label{eq:exact-secant-linear-equation}
\end{equation}
Estimate \eqref{eq:quantitative-directional-coarea} applies uniformly
to every convex combination \((1-t)U+tU_h\).  Together with Lemma
\ref{lem:regional-angular-compactness}, it gives, for each compact
\(K\Subset B_{r_0}\), a number \(h_K>0\) such that
\begin{equation}
 \sup_{0<h<h_K}
 \int_K\int_{\R^n}(1\wedge|x-y|^2)
 \frac{c_h(x,y)}{|x-y|^{n+sp}}\dd y\!\dd x<\infty.
 \label{eq:averaged-secant-levy-bound}
\end{equation}
after reducing the admissible upper bound for \(h\), if necessary.
For clarity, on the near-diagonal part the angular lemma gives
\[
 c_h(x,x-\rho\omega)=\rho^{p-2}b_h(x,\rho,\omega),
 \qquad
 \int_{\mathbb S^{n-1}}b_h(x,\rho,\omega)\dd\omega\le C,
\]
and the corresponding radial integral is
\[
 \int_0^1\rho^{p-2}\rho^2\rho^{-sp-1}\dd\rho
 =\int_0^1\rho^{1-\tau}\dd\rho<\infty.
\]
On separated sets,
\eqref{eq:quantitative-directional-coarea} gives uniform
integrability.  More precisely, if \(K\Subset B_{r_0}\),
\(0<\varepsilon<R<\infty\), and
\[
 G\subset\{(x,y):x\in K,\ \varepsilon\le|x-y|\le R\},
\]
then application of that estimate to each \(x\)-fiber, followed by
H\"older's inequality in \(y\), gives
\begin{equation}
 \sup_{0<h<h_K}\iint_G
 \frac{c_h(x,y)}{|x-y|^{n+sp}}\dd x\!\dd y
 \le C_{K,\varepsilon,R}|G|^{p-1}.
 \label{eq:separated-secant-uniform-integrability}
\end{equation}

Choose a cutoff \(\eta\in C_c^\infty(B_{r_0})\) and test
\eqref{eq:exact-secant-linear-equation} with $\eta^2w_h$.  Since
\(|w_h|\le L\), the identity
\[
 (a-b)(\eta_x^2a-\eta_y^2b)
 =|\eta_xa-\eta_yb|^2-ab|\eta_x-\eta_y|^2
\]
and \eqref{eq:averaged-secant-levy-bound} give
\begin{equation}
 \sup_{0<h<h_\eta}
 \iint_{\R^n\times\R^n}
 \frac{c_h(x,y)|\delta(\eta w_h)(x,y)|^2}
 {|x-y|^{n+sp}}\dd x\!\dd y<\infty.
 \label{eq:uniform-secant-caccioppoli}
\end{equation}
where \(h_\eta\) is the number attached above to a compact
neighborhood of \(\operatorname{supp}\eta\).  We clarify the
admissibility of this test without invoking density for the measurable
secant form.  For each fixed \(h>0\),
\(U,U_h\in W^{s,p}_{\rm loc}\), and
\[
 c_h|\delta w_h|^2
 =h^{-2}\bigl[\Jp(\delta U_h)-\Jp(\delta U)\bigr]
             (\delta U_h-\delta U)
 \le Ch^{-2}|\delta U_h-\delta U|^p.
\]
Moreover, \(\eta^2w_h=\eta^2(U_h-U)/h\in W^{s,p}_0(B_{r_0})\).
Approximate this function in \(W^{s,p}_0(B_{r_0})\) by smooth compactly
supported functions, use the same approximants separately in the weak
equations for \(U_h\) and \(U\), and then subtract.  The local
\(W^{s,p}\) bounds and the global Lipschitz tail, which is integrable
because \(sp>p-1\), justify the limit in the two original weak forms.
The secant identity then gives
\eqref{eq:exact-secant-linear-equation} with the test
\(\eta^2w_h\).  The restrictions
\(|x-y|>\varepsilon\) and \(c_h\le N\) are used only to estimate the
resulting integrals; they do not replace the kernel in the equation.
Letting first \(N\uparrow\infty\) and then
\(\varepsilon\downarrow0\), the averaged L\'evy bound and Fatou's
lemma give \eqref{eq:uniform-secant-caccioppoli}.

As $h\downarrow0$, $w_h\to w$ locally uniformly in $B_{r_0}$ and almost
everywhere in $\R^n$; the latter follows from Rademacher's theorem, or
equivalently from one-dimensional differentiation on almost every line
parallel to \(e\).  For almost every pair with one point in
$B_{r_0}$,
\[
 \frac{c_h(x,y)}{|x-y|^{n+sp}}\longrightarrow J_U(x,y).
\]
Indeed, the exceptional set \(U(x)=U(y)\) has zero measure in \(x\)
for each fixed \(y\), by strict monotonicity on the \(e\)-lines.
Fatou's lemma in \eqref{eq:uniform-secant-caccioppoli} first gives the
local \(J_U\)-energy of \(w\).

We verify at this point that the limiting energy belongs to the local
closed form domain, rather than only to the maximal finite-energy
class.  We first prove the required core property directly on
\(D_0:=B_{r_0}\).  Define
\[
 \Theta(x):=P_{e^\perp}x+U(x)e,
 \qquad x\in D_0.
\]
The strict bound \(\partial_eU\ge\lambda_0\), the local Lipschitz
bound, and integration on the \(e\)-lines show that \(\Theta\) is
bi-Lipschitz from \(D_0\) onto \(D^\sharp:=\Theta(D_0)\); moreover,
\(\det D\Theta=\partial_eU\) is bounded above and below.  Indeed,
\(\Theta\) preserves the transverse coordinates, while strict
monotonicity on each \(e\)-line controls the remaining coordinate of
\(\Theta^{-1}\).  More explicitly, after rotating so that \(e=e_n\)
and inserting the intermediate point \((x',y_n)\in B_{2r_0}\), one has
\[
 |x_n-y_n|\le\lambda_0^{-1}
 \bigl(|U(x)-U(y)|+L|x'-y'|\bigr),
\]
which gives \(|x-y|\le C|\Theta(x)-\Theta(y)|\).

Let \(f\in L^2(D_0)\) have compact support in \(D_0\) and finite
regional \(J_U\)-energy, and put
\(\widetilde f=f\circ\Theta^{-1}\).  With
\(\xi=\Theta(x)\) and \(\zeta=\Theta(y)\), change of variables gives
\begin{equation}
 \iint_{D_0\times D_0}J_U(x,y)|\delta f(x,y)|^2\dd x\!\dd y
 \asymp
 \iint_{D^\sharp\times D^\sharp}
 \frac{|e\cdot(\xi-\zeta)|^{p-2}
       |\delta\widetilde f(\xi,\zeta)|^2}
 {|\xi-\zeta|^{n+sp}}\dd\xi\!\dd\zeta.
 \label{eq:noncritical-coordinate-form-equivalence}
\end{equation}
The kernel on the right is the translation-invariant kernel
\[
 L(H):=\frac{|e\cdot H|^{p-2}}{|H|^{n+sp}}
 =\frac{|e\cdot\omega|^{p-2}}{|H|^{n+\tau}},
 \qquad \omega=\frac H{|H|}.
\]
We denote its full-space quadratic form by
\[
 \mathcal E_L(g):=\iint_{\R^n\times\R^n}
 L(\xi-\zeta)|g(\xi)-g(\zeta)|^2\dd\xi\!\dd\zeta.
\]
Its angular density is integrable because \(p-2>-1\), and its Fourier
symbol satisfies
\[
 c_{n,\tau}\int_{\mathbb S^{n-1}}
 |\vartheta\cdot\omega|^\tau|e\cdot\omega|^{p-2}\dd\omega
 \asymp |\vartheta|^\tau.
\]
Thus its full-space form domain is \(H^{\tau/2}(\R^n)\).

Extend \(\widetilde f\) by zero outside \(D^\sharp\).  Since its
support has positive distance from \(\partial D^\sharp\),
\eqref{eq:noncritical-coordinate-form-equivalence} and direct radial
integration show that its full-space \(L\)-energy is finite.  If
\(\rho_\varepsilon\) is a standard mollifier, translation invariance
and Jensen's inequality give
\[
 \mathcal E_L(\rho_\varepsilon*\widetilde f-\widetilde f)
 \le\int_{\R^n}\rho_\varepsilon(h)
       \mathcal E_L(\widetilde f(\,\cdot+h)-\widetilde f)\dd h
 \longrightarrow0.
\]
For completeness, the last convergence follows by writing the energy
as the integral in \(H\) of the squared \(L^2\)-norm of the
corresponding increment.  Translation continuity in \(L^2\) gives
pointwise convergence in \(H\), while
\(4L(H)\|\widetilde f(\,\cdot+H)-\widetilde f\|_2^2\) is an integrable
majorant.  For small \(\varepsilon\), the convolution remains compactly
supported in \(D^\sharp\).

Pulling it back by \(\Theta\) gives a compactly supported
\(C^{1,\alpha}\) function in \(D_0\), converging to \(f\) in
\(L^2\) and in the regional \(J_U\)-energy by
\eqref{eq:noncritical-coordinate-form-equivalence}.  Mollifying once
more in the original variables produces functions in
\(C_c^\infty(D_0)\).  For this second mollification the differences
converge in \(C^1\); Lemma~\ref{lem:regional-angular-compactness} and
\(\int_0^1r^{1-\tau}\dd r<\infty\) then give convergence in the
\(J_U\)-energy.  If \(f\ge0\), both mollifications preserve
nonnegativity.  We have
therefore proved that every compactly supported finite-energy function
in \(D_0\) belongs to \(\mathcal F_{J_U}(D_0)\).  The approximants may
be chosen with supports in one fixed compact subset of \(D_0\).

Fatou's lemma applied to \eqref{eq:uniform-secant-caccioppoli} gives
finite \(J_U\)-energy for \(\eta w\), not merely for \(w\) on pairs
inside the cutoff set.  Applying the preceding core property to
\(f=\eta w\) proves
\(\eta w\in\mathcal F_{J_U}(B_{r_0})\).  Thus
\(w\in\mathcal F_{J_U,\mathrm{loc}}(B_{r_0})\).

To pass to the equation, put
\(K=\operatorname{supp}\varphi\) and choose the preceding cutoff
\(\eta\) to be one in a fixed neighborhood of \(K\).  We spell out the
three truncations used in \eqref{eq:exact-secant-linear-equation}.

First, on
\[
 \{(x,y):x\in K,
       \varepsilon\le|x-y|\le R\},
\]
the uniform integrability following
\eqref{eq:quantitative-directional-coarea}, the almost everywhere
convergence of the kernels, and \(|w_h|\le L\) permit the use of
Vitali's theorem.  Hence the integral converges to the corresponding
integral with kernel \(J_U\).

Second, if \(|y|>R\), then \(\delta\varphi(x,y)=\varphi(x)\) and
\(|\delta w_h(x,y)|\le2L\).  Directional coarea in the \(x\)-variable
gives
\begin{equation*}
 \sup_{\substack{0<h<h_K\\t\in[0,1]}}\int_K
 |(1-t)\delta U+t\delta U_h|^{p-2}\dd x\le C
\end{equation*}
for each \(y\).  After multiplication by
\(|x-y|^{-n-sp}\) and integration in \(y\), the contribution is
\(O(R^{-sp})\), uniformly in \(h\).

Third, choose \(\varepsilon\) smaller than the distance from \(K\) to
\(\{\eta\ne1\}\).  Then every pair with \(|x-y|<\varepsilon\) which
meets \(K\) lies in \(\{\eta=1\}\).  Cauchy--Schwarz,
\eqref{eq:uniform-secant-caccioppoli}, and the angular upper bound
therefore give
\[
 C\left(\int_0^\varepsilon r^{1-\tau}\dd r\right)^{1/2}
 =C\varepsilon^{(2-\tau)/2}\longrightarrow0.
\]
The estimate is uniform in \(h\).  Letting successively
\(h\downarrow0\), \(R\uparrow\infty\), and
\(\varepsilon\downarrow0\) permits passage to the limit in
\eqref{eq:exact-secant-linear-equation}.  This proves
\eqref{eq:full-weak-linearization}.

It remains to justify the asserted extension of the identity.  Let
\(\varphi\in\mathcal F_{J_U,c}(B_{r_0})\).  Multiplying a form-norm
approximating sequence by a fixed cutoff which equals one near
\(\operatorname{supp}\varphi\), we may choose
\(\varphi_j\in C_c^\infty(B_{r_0})\) with supports in a common compact
set and
\[
 \|\varphi_j-\varphi\|_2^2
 +\iint_{B_{r_0}\times B_{r_0}}J_U
       |\delta(\varphi_j-\varphi)|^2\dd x\!\dd y\longrightarrow0.
\]
Choose \(K\Subset K_1\Subset B_{r_0}\) so that all the supports lie in
\(K\).  On \(K_1\times K_1\), the regional part of
\eqref{eq:full-weak-linearization} passes to the limit by
Cauchy--Schwarz and the local form membership of \(w\).  Pairs meeting
\(K\) and \(B_{r_0}\setminus K_1\) are separated; directional coarea
gives the same local Morrey bound for their coefficient, and
\eqref{eq:kato-infinitesimal-form-bound} makes their contribution tend
to zero.  For the crossed part, write
\(A_U(x)=\int_{\R^n\setminus B_{r_0}}
J_U(x,y)\dd y\).  Since \(|w|\le L\),
\begin{align*}
 &\left|\int_{B_{r_0}}\int_{\R^n\setminus B_{r_0}}
 J_U(x,y)\delta w(x,y)(\varphi_j-\varphi)(x)\dd y\!\dd x\right|\\
 &\qquad\le2L
 \left(\int_KA_U\dd x\right)^{1/2}
 \left(\int_KA_U|\varphi_j-\varphi|^2\dd x\right)^{1/2}.
\end{align*}
The same
directional-coarea argument used for \eqref{eq:crossed-morrey-bound},
now with constants depending also on
\(\operatorname{dist}(K,\partial B_{r_0})\), gives the required local
Morrey bound for \(A_U\) on a neighborhood of \(K\).  The last factor
therefore tends to zero by the density extension of
\eqref{eq:kato-infinitesimal-form-bound}, applied to the kernel
\(J_U\).  Hence
\eqref{eq:full-weak-linearization} holds for every compactly supported
form-domain test.
\end{proof}

\end{document}
